\subsection{连接振荡电路并观察输出信号}

将振荡电路连接到电压源，并观察信号随电容器浸入深度变化而发生的变化。对于处于空气中的电容器，测得信号频率为 $\SI{948.35}{\kilo\hertz}$。

\subsection{连接比较器}

将振荡电路的输出与比较器连接，以便将信号数字化。调整比较器，使其在某一电压阈值处切换，从而产生矩形输出信号。由此可以精确确定信号频率，并尽量减小可能的干扰。

\subsection{读取输出信号频率}

第一组测量中，将电容器完全浸入液体。随后使用频率测量仪记录由此得到的电容以及输出频率。为了获得可靠且可重复的测量值，该过程重复多次。

由图可得：$f(\SI{23}{\centi\meter}) = \SI{343.371}{\kilo\hertz}$。

随后，将电容器在多个规定液位下部分浸入液体。对于每个规定液位，测量并记录相应的电容值和输出频率。

由图可得：$f(\SI{34}{\centi\meter}) = \SI{372.657}{\kilo\hertz}$。

\subsection{使用测量卡读取信号值}

振荡电路的信号值通过测量卡传输到计算机，随后在 MATLAB 中进行分析。为此使用实验计算机桌面上 \texttt{Matlab\_Vorlage} 文件夹中的脚本 \texttt{main.m}。必须为测量信号设置合适的采样率。

（采样的真实数据，这图是参考报告里的示例）

\subsection{Schmitt 触发器函数}

当达到正向开关阈值 $U_{\mathrm{ein}}$ 时，Schmitt 触发器将输出电压切换到最大值 $U_{a,\max}$。一旦超过负向开关阈值 $U_{\mathrm{aus}}$，它就切换到最小值 $U_{a,\min}$。在 MATLAB 中，该函数实现如下：

（施密特代码截图，这图是参考报告里的示例）
（画出的波形图，这图是参考报告里的示例）
从图中可以看出，$T = \_$，$N_x = \_$，$f_x = N_x/T = \_$。
（画出的波形图，这图是参考报告里的示例）

\subsection{步骤自动化}

编写的工作流程通过循环实现自动化。由此可以连续计算频率值，并直接在屏幕上显示。这样无需进一步手动干预即可持续监测和分析信号，从而保证测量结果准确且一致。

（自动化代码截图，这图是参考报告里的示例）
